%% Dans ce fichier, les �ventuels accents sont cod�s
%% en ISO-8859-15 (aka latin9)
%% La ligne suivante donne trois exemples 
%% e accent aigu : �, e accent grave : �, e-dans-l'o : �
%%%%%ENTREES%%%%
\begin{tabarg}
\multicolumn{3}{l}{\titretab {Entr\'ees}}\\
\hline
$\rho_a$ & Vecteur, \texttt{float} & Densit� de l'air sec ($kg.m^{-3}$)\\
$\rho_v$ & Vecteur, \texttt{float} & Concentration en vapeur d'eau ($g/m^3$)\\
$\rho_c$ & Vecteur, \texttt{float} & Concentration du compos� dont on veut corriger le flux ($g/m^3$)\\
$T_a$ & Vecteur, \texttt{float} & Temp�rature de l'air ($^\circ$ C ou K)\\
$Cov\_w\rho_v$ & Vecteur, \texttt{float} & Covariance $\overline{w'\rho_v'}$ ($g.s^{-1}.m^{-2}$)\\
$Cov\_wt$ & Vecteur, \texttt{float} & Covariance $\overline{w'T_a'}$ (m.K.s-1)\\
$Cov\_w\rho_c$ & Vecteur, \texttt{float} & Covariance $\overline{w'\rho_c'}$ sur laquelle on va appliquer la correction ($g.s^{-1}.m^{-2}$)\\
\end{tabarg}

\medskip
%%%%%SORTIES%%%%
\begin{tabarg}
\multicolumn{3}{l}{\titretab {Sorties}}\\
\hline
$w_d$ & Vecteur, \texttt{float} & Vitesse verticale moyenne d�e au flux de densit� (m/s)\\
$Flx_{WPL}$ & Vecteur, \texttt{float} &  Flux du compos� corrig� du flux de densit� selon WPL ($g.s^{-1}.m^{-2}$)\\
\end{tabarg}
